\documentclass[11pt]{article}

% ---------------------------------------------------------
% Packages
% ---------------------------------------------------------
\usepackage[a4paper, margin=1in]{geometry}
\usepackage{amsmath, amssymb, amsthm}
\usepackage{setspace}
\usepackage{graphicx}
\usepackage{float}
\usepackage{caption}
\usepackage{subcaption}
\usepackage{booktabs}
\usepackage{hyperref}
\usepackage{enumitem}
\usepackage{mathtools}
\usepackage{physics}
\usepackage{bm}
\usepackage{listings}
\usepackage{xcolor}


% ---------------------------------------------------------
% Code block style
% ---------------------------------------------------------
\lstset{
    basicstyle=\ttfamily\small,
    keywordstyle=\color{blue},
    commentstyle=\color{gray},
    stringstyle=\color{orange},
    breaklines=true,
    columns=fullflexible
}

% ---------------------------------------------------------
% indentation
% ---------------------------------------------------------
\setlength{\parindent}{0pt}


% ---------------------------------------------------------
% Hyperlink setup
% ---------------------------------------------------------
\hypersetup{
    colorlinks=true,
    linkcolor=black,
    urlcolor=blue,
    citecolor=blue
}

% ---------------------------------------------------------
% Title
% ---------------------------------------------------------
\title{
    \vspace{-1cm}
    \textbf{Lab 3 Report}\\
    \vspace{0.2cm}
    Summer 2026
}
\author{
    Soeun Jang \\
    Virginia Tech \\
    \texttt{sophiaj04@vt.edu}
}
\date{\today}

% ---------------------------------------------------------
% Document
% ---------------------------------------------------------
\begin{document}

\maketitle
\newpage   

\tableofcontents
\newpage   

% ---------------------------------------------------------
\section*{Introduction}
\addcontentsline{toc}{section}{Introduction}
% ---------------------------------------------------------



% ---------------------------------------------------------
\section*{Part 1: Tensor Basics}
\addcontentsline{toc}{section}{Part 1: Tensor Basics}
% ---------------------------------------------------------

\subsection*{Exercise 1}
\addcontentsline{toc}{subsection}{Exercise 1}
\textbf{1}. Unfold function

\begin{lstlisting}
function unfold(A::Array{Float64,3}, n::Int)::Matrix{Float64}
    n1, n2, n3 = size(A)
    if n == 1
        return reshape(permutedims(A, [1,2,3]), n1, n2*n3)
    elseif n == 2
        return reshape(permutedims(A, [2,1,3]), n2, n1*n3)
    else
        return reshape(permutedims(A, [3,1,2]), n3, n1*n2)
    end
end

\end{lstlisting}

Based on the function given, we implemented unfold function that returns those values. To verify this function works, 

\begin{lstlisting}
function unfold(A::Array{Float64,3}, n::Int)::Matrix{Float64}
    n1, n2, n3 = size(A)
    if n == 1
        return reshape(permutedims(A, [1,2,3]), n1, n2*n3)
    elseif n == 2
        return reshape(permutedims(A, [2,1,3]), n2, n1*n3)
    else
        return reshape(permutedims(A, [3,1,2]), n3, n1*n2)
    end
end

\end{lstlisting}

We used this verification and every assertion passed which confirms that the implemented unfold function successfully rearranges the tenor structure into appropriate 2-dimensional matrix form for each mode.

\vspace{1.5em}

\textbf{2}. Implement mode product function

\begin{lstlisting}
function mode_product(A::Array{Float64,3}, M::Matrix{Float64}, n::Int)::Array{Float64,3}
    n1, n2, n3 = size(A)
    s = size(M, 1)
    
    An = unfold(A, n)
    B_unfolded = M * An
    
    if n == 1
        return permutedims(reshape(B_unfolded, s, n2, n3), [1,2,3])
    elseif n == 2
        return permutedims(reshape(B_unfolded, s, n1, n3), [2,1,3])
    else
        return permutedims(reshape(B_unfolded, s, n1, n2), [2,3,1])
    end
end
\end{lstlisting}

The mode n product between a 3 way tensor and a matrix, those are defined as multiplying M with the mode unfolding of A then just folding the result back into the size tensor.
To verify this function,

\begin{lstlisting}
#Ex 1-2
A = rand(4, 5, 6)
M1 = rand(3, 4)
B = mode_product(A, M1, 1)
println("mode_product(A, M1, 1) size: ", size(B))
@assert size(B) == (3, 5, 6)
println("passed for mode_product with n=1.")
\end{lstlisting}

This prints out
"mode\_product(A, M1, 1) size: (3, 5, 6)
passed for mode\_product with n=1." which indicates that it passes. The resulting tensor has a size of (3,5,6) which means that it confirms the implementation is correct since multiplication of mode-1 changes into 3. 


\vspace{1.5em}

\textbf{3}. Verification of Commutativity error

\begin{lstlisting}
#Ex 1-3
M2 = rand(7, 5)

C1 = mode_product(mode_product(A, M1, 1), M2, 2)
C2 = mode_product(mode_product(A, M2, 2), M1, 1)

err = maximum(abs.(C1 - C2))
println("commutativity error: ", err)
@assert err < 1e-12
println("passed for commutativity of mode_product.")
\end{lstlisting}

Based on same A, M1, commutativity error: 8.881784197001252e-16
passed for commutativity of mode\_product. Since the multiplication of mode-1 and mode-2 are commuting each other which means they are always identical tensor.

\vspace{1.5em}

\textbf{4}. Separable tensor verification

\begin{lstlisting}
n = 20
pts = range(0, pi, length=n)

f1 = sin.(pts)
f2 = cos.(2 .* pts)
f3 = exp.(.-pts)

# A[i,j,k] = f1(xi) * f2(yj) * f3(zk)
A_sep = [f1[i] * f2[j] * f3[k] for i in 1:n, j in 1:n, k in 1:n]

A1 = unfold(A_sep, 1)
A2 = unfold(A_sep, 2)
A3 = unfold(A_sep, 3)

sv1 = svdvals(A1)
sv2 = svdvals(A2)
sv3 = svdvals(A3)

tol = 1e-10
r1 = sum(sv1 .> tol)
r2 = sum(sv2 .> tol)
r3 = sum(sv3 .> tol)

println("rank of mode-1 unfolding: ", r1)
println("rank of mode-2 unfolding: ", r2)
println("rank of mode-3 unfolding: ", r3)
\end{lstlisting}

This prints out rank of mode-1 unfolding: 1
rank of mode-2 unfolding: 1
rank of mode-3 unfolding: 1 since every separated tensor $A_{ijk} = f_1(x_i)f_2(yIj)f_3(z_k)$ contains three vector's external structure, three mode unfolding rank should be all 1. 

% ---------------------------------------------------------
\section*{Part 2:  HOSVD and the Tucker Struct}
\addcontentsline{toc}{section}{Part 2:  HOSVD and the Tucker Struct}
% ---------------------------------------------------------
\subsection*{Exercise 2}
\addcontentsline{toc}{subsection}{Exercise 2}
\textbf{1}. Define the Tucker data structure

\begin{lstlisting}
mutable struct Tucker3
    G::Array{Float64,3}  # r1 x r2 x r3 core tensor
    U1::Matrix{Float64}  # n1 x r1
    U2::Matrix{Float64}  # n2 x r2
    U3::Matrix{Float64}  # n3 x r3
end
\end{lstlisting}

This tucker3 structure is the data structure to store tucker decomposition by G which is small core tensor and U1, U2, U3 works as a base that is the factor matrices of the mode. 
Therefore, this structure captures both the compressed core and the mode wise basis matrices that reconstruct the full tensor.

\vspace{1.5em}

\textbf{2}. Implement todense 

\begin{lstlisting}
function todense(T::Tucker3)::Array{Float64,3}
    A = mode_product(T.G, T.U1, 1)
    A = mode_product(A, T.U2, 2)
    A = mode_product(A, T.U3, 3)
    return A
end
\end{lstlisting}

The todense function reconstructs the full tensor by applying the tree mode products and it expands the compressed core G tensor back to its original dimensions.

\vspace{1.5em}

\textbf{3}. Implement hosvd function

\begin{lstlisting}
function hosvd(A::Array{Float64,3}, r::Tuple{Int,Int,Int})::Tucker3
    r1, r2, r3 = r

    U1 = svd(unfold(A, 1)).U[:, 1:r1]
    U2 = svd(unfold(A, 2)).U[:, 1:r2]
    U3 = svd(unfold(A, 3)).U[:, 1:r3]

    G = mode_product(A, U1', 1)
    G = mode_product(G, U2', 2)
    G = mode_product(G, U3', 3)

    return Tucker3(G, U1, U2, U3)
end
\end{lstlisting}

In the HOSVD procedure, we first compute the mode‑1, mode‑2, and mode‑3 unfoldings of the tensor A. From each unfolding, we apply SVD and take the leading r1, r2, r3 to form the factor matrices U1, U2, U3 each. After obtaining these factor matrices, we compute the core tensor G by multiplying the original tensor A with the transposes of the factor matrices along each side and this core tensor G with U1, U2, U3 forms the tucker3 decomposition.

\vspace{1.5em}

\textbf{4}. Implement Tucker3 function

\begin{lstlisting}
function Tucker3(A::Array{Float64,3}, tol::Float64)::Tucker3
    threshold = tol / sqrt(3) * norm(A)

    F1 = svd(unfold(A, 1))
    F2 = svd(unfold(A, 2))
    F3 = svd(unfold(A, 3))

    r1 = max(1, sum(F1.S .>= threshold))
    r2 = max(1, sum(F2.S .>= threshold))
    r3 = max(1, sum(F3.S .>= threshold))

    U1 = F1.U[:, 1:r1]
    U2 = F2.U[:, 1:r2]
    U3 = F3.U[:, 1:r3]

    # Core tensor
    G = mode_product(A, Matrix(U1'), 1)
    G = mode_product(G, Matrix(U2'), 2)
    G = mode_product(G, Matrix(U3'), 3)

    return Tucker3(G, U1, U2, U3)
end
\end{lstlisting}

This function makes tensor A be compressed as the way of tucker. Based on the tolerance called tol, it chooses the multi-linear ranks r value that is greater than threshold. Then, it applies SVD for each mode unfolding A to get the factor matrices U1, U2, U3 that correspond to each directions of modes 1,2,3. To get core tensor G by compressing each directions and return the structure with tucker3 structure. 

\vspace{1.5em}

\textbf{5}. Verification

\begin{lstlisting}
include("../src/unfold.jl")
include("../src/mode.product.jl")
include("../src/tucker3.jl")
include("../src/tucker_3.jl")
include("../src/todense.jl")
include("../src/hosvd.jl")


using LinearAlgebra

A = randn(30, 30, 30)
T = Tucker3(A, 1e-12)

err = norm(todense(T) - A)
println("norm of A - todense(Tucker3(A, 1e-12))= ", err)
println("1e-10 * norm(A)= ", 1e-10 * norm(A))
@assert err < 1e-10 * norm(A)
println("Verified that the error is within the specified tolerance.")
\end{lstlisting}

Based on this verification, it prints out "norm of A - todense(Tucker3(A, 1e-12))= 3.8296095413294746e-13" and "1e-10 * norm(A)= 1.6477528417710177e-8". It verifies that the error is within the specified tolerance. In conclusion, tucker3 decomposition works correctly.

\section*{Part 3: Tucker Truncation and Error Analysis}
\addcontentsline{toc}{section}{Part 3: Tucker Truncation and Error Analysis}
\subsection*{Exercise 3}
\addcontentsline{toc}{subsection}{Exercise 3}
\textbf{1}. Implement tucker round function
\begin{lstlisting}
function tucker_round(T::Tucker3, tol::Float64)::Tucker3
    G  = T.G
    U1 = T.U1
    U2 = T.U2
    U3 = T.U3

    F1 = svd(unfold(G, 1))
    F2 = svd(unfold(G, 2))
    F3 = svd(unfold(G, 3))

    threshold = tol / sqrt(3) * norm(G)

    r1 = max(1, sum(F1.S .>= threshold))
    r2 = max(1, sum(F2.S .>= threshold))
    r3 = max(1, sum(F3.S .>= threshold))

    Ub1 = F1.U[:, 1:r1]
    Ub2 = F2.U[:, 1:r2]
    Ub3 = F3.U[:, 1:r3]

    Gnew = mode_product(G, Matrix(Ub1'), 1)
    Gnew = mode_product(Gnew, Matrix(Ub2'), 2)
    Gnew = mode_product(Gnew, Matrix(Ub3'), 3)

    U1new = U1 * Ub1
    U2new = U2 * Ub2
    U3new = U3 * Ub3

    return Tucker3(Gnew, U1new, U2new, U3new)
end

\end{lstlisting}

This function explains the tucker round which applies HOSVD to the core G, then update the factor matrices without ever expanding to the full tensor.

\vspace{1.5em}

\textbf{2}. Test

\begin{lstlisting}
xvinclude("../src/unfold.jl")
include("../src/mode.product.jl")
include("../src/tucker3.jl")     
include("../src/tucker_3.jl")   
include("../src/todense.jl")
include("../src/hosvd.jl")
include("../src/tucker_round.jl") 

using Test
using LinearAlgebra

@testset "Tucker_Round Test" begin
    n1 = n2 = n3 = 30
    r1 = r2 = r3 = 20

    U1 = randn(n1, r1)
    U2 = randn(n2, r2)
    U3 = randn(n3, r3)
    G  = randn(r1, r2, r3)

    A = mode_product(mode_product(mode_product(G, U1, 1), U2, 2), U3, 3)

    T = Tucker3(A, 1e-12)

    tols = [1e-2, 1e-4, 1e-8]

    for tol in tols
        Tr = tucker_round(T, tol)

        r = (size(Tr.G,1), size(Tr.G,2), size(Tr.G,3))

        err = norm(todense(Tr) - A) / norm(A)

        @test err < 10 * tol  
        @test all(ri <= 20 for ri in r)

        println("tol = $tol   ranks = $r   rel error = $err")
    end
end

\end{lstlisting}

This test gives the error value about $3.41e-15$ since we are rounding the core which is already orthogonally compressed for every tolerance and rank $(20,20,20)$. Since the synthetic tensor actually has true multi-linear rank, all tolerances give rank of $20$. Also the singular values of the core are all large, and non fall below $10^{-2}, 10^{-4}, 10^{-8}$ which means there is no truncation so the rounded ranks stay at $20,20,20$. 

\vspace{1.5em}

\textbf{3}. HOSVD multi-linear rank

\begin{lstlisting}
@testset "Gaussian Tucker Rank Test" begin
    n = 100
    xs = range(-3, 3, length=n)
    ys = range(-3, 3, length=n)
    zs = range(-3, 3, length=n)

    A = Array{Float64,3}(undef, n, n, n)
    for i in 1:n, j in 1:n, k in 1:n
        A[i,j,k] = exp(-(xs[i]^2 + ys[j]^2 + zs[k]^2))
    end

    T = Tucker3(A, 1e-12)

    r = (size(T.G,1), size(T.G,2), size(T.G,3))
    println("Multilinear rank assigned by HOSVD = ", r)

    @test r[1] >= 1
    @test r[2] >= 1
    @test r[3] >= 1

    err = norm(todense(T) - A) / norm(A)
    println("Error = ", err)

    @test err < 1e-10
end
\end{lstlisting}

Confirming 3D Gaussian Tucker rank, this returns the multilinear rank (1,1,1) with error about 1.12e-14. HOSVD detects the correct rank. Also the error is tiny, which is an almost exact representation of the discrete tensor.

\subsection*{Exercise 4}
\addcontentsline{toc}{subsection}{Exercise 4}
\begin{lstlisting}
using Test
using LinearAlgebra
using Plots

# helper function to unfold a 3D array along mode n
function unfold(A, n)
    if n == 1
        return reshape(A, size(A,1), :)
    elseif n == 2
        return reshape(permutedims(A, (2,1,3)), size(A,2), :)
    else
        return reshape(permutedims(A, (3,1,2)), size(A,3), :)
    end
end

# singular values of mode-n unfolding
function mode_svals(A)
    s1 = svd(unfold(A,1)).S
    s2 = svd(unfold(A,2)).S
    s3 = svd(unfold(A,3)).S
    return s1, s2, s3
end

@testset "Exercise 4: Singular Value Decay" begin
    n = 50

    #1
    xs = range(0, 2π, length=n)
    ys = range(0, 2π, length=n)
    zs = range(0, 2π, length=n)

    A1 = [sin(x+y+z) for x in xs, y in ys, z in zs]
    s1a, s2a, s3a = mode_svals(A1)

    r1a = sum(s1a .> 1e-6)
    r2a = sum(s2a .> 1e-6)
    r3a = sum(s3a .> 1e-6)

    println("\n(1) sin(x+y+z) ranks = ", (r1a, r2a, r3a))

    #2
    eps = 0.1
    xs = range(0, 1, length=n)
    ys = range(0, 1, length=n)
    zs = range(0, 1, length=n)

    A2 = [1/sqrt(x^2 + y^2 + z^2 + eps^2) for x in xs, y in ys, z in zs]
    s1b, s2b, s3b = mode_svals(A2)

    r1b = sum(s1b .> 1e-6)
    r2b = sum(s2b .> 1e-6)
    r3b = sum(s3b .> 1e-6)

    println("(2) 1/sqrt(x^2+y^2+z^2+eps^2) ranks = ", (r1b, r2b, r3b))

    #3
    xs = range(-1, 1, length=n)
    ys = range(-1, 1, length=n)
    zs = range(-1, 1, length=n)

    A3 = [tanh(10*(x^2 + y^2 - z)) for x in xs, y in ys, z in zs]
    s1c, s2c, s3c = mode_svals(A3)

    r1c = sum(s1c .> 1e-6)
    r2c = sum(s2c .> 1e-6)
    r3c = sum(s3c .> 1e-6)

    println("(3) tanh(10(x^2+y^2−z)) ranks = ", (r1c, r2c, r3c))

    #plotting singular value decay
    p1 = plot(s1a, yscale=:log10, label="mode-1", title="sin(x+y+z)")
    plot!(p1, s2a, label="mode-2")
    plot!(p1, s3a, label="mode-3")
    savefig(p1, "p1.png")
    display(p1)

    p2 = plot(s1b, yscale=:log10, label="mode-1", title="1/sqrt(x^2+y^2+z^2+eps^2)")
    plot!(p2, s2b, label="mode-2")
    plot!(p2, s3b, label="mode-3")
    savefig(p2, "p2.png")
    display(p2)

    p3 = plot(s1c, yscale=:log10, label="mode-1", title="tanh(10(x^2+y^2-z))")
    plot!(p3, s2c, label="mode-2")
    plot!(p3, s3c, label="mode-3")

    savefig(p3, "p3.png")
    display(p3)

    plot(p1, p2, p3, layout=(3,1), size=(800,1200))
end

\end{lstlisting}

\textbf{1}. $sin(x + y + z)$

\vspace{1em}

The ranks are (2, 2, 2) which is symmetric balanced. Each mode unfolding sees the same structure and the singular values decay identically across modes. This is the most balanced case. 

\vspace{1.5em}

\textbf{2}. $ 1/\sqrt{x^2+y^2+z^2+\epsilon^2} $

\vspace{1em}

The ranks are (14, 14, 14) which is balanced but it contains higher rank. This function is radially symmetric which means the function depends only on the distance from the origin.

\vspace{1.5em}

\textbf{3}. $ tanh(10(x^2+y^2-z))$

\vspace{1em}

The ranks are (24, 24, 29) which is unbalanced. This function treats the variables asymmetrically. $x,y$ shows quadratically and symmetrically but $z$ appears linearly and contains opposite sign. 

\section*{Part 4: Tucker Arithmetic and 3D PDE Setup}
\addcontentsline{toc}{section}{Part 4: Tucker Arithmetic and 3D PDE Setup}
\subsection*{Exercise 5}
\addcontentsline{toc}{subsection}{Exercise 5}
\textbf{1}.

\vspace{1.5em}

\textbf{2}.

\vspace{1.5em}

\textbf{3}.

\subsection*{Exercise 6}
\addcontentsline{toc}{subsection}{Exercise 6}
\textbf{1}.

\vspace{1.5em}

\textbf{2}.

\vspace{1.5em}

\textbf{3}.

\vspace{1.5em}

\textbf{4}.


\subsection*{Exercise 7}
\addcontentsline{toc}{subsection}{Exercise 7}

\textbf{1}.

\vspace{1.5em}

\textbf{2}.

\vspace{1.5em}

\textbf{3}.

\vspace{1.5em}

\textbf{4}.

\vspace{1.5em}

\textbf{5}.


% ---------------------------------------------------------
\section*{Conclusion}
\addcontentsline{toc}{section}
{Conclusion}





\vspace{2em}

 

% ---------------------------------------------------------


\end{document}
